\thispagestyle{empty}

\begin{center}
    \fontsize{12pt}{12pt}\selectfont
    \begin{minipage}{0.15\textwidth}
        \includegraphics[width=\linewidth]{img/bmstu_logo.jpg}
    \end{minipage}
    \hfill
    \begin{minipage}{0.80\textwidth}\centering\bfseries
        {\linespread{1}\selectfont
        {Министерство науки и высшего образования Российской Федерации}

        {Федеральное государственное бюджетное~образовательное~учреждение высшего образования}

        {<<Московский государственный технический университет имени~Н.~Э.~Баумана (национальный~исследовательский~университет)>>}

        {(МГТУ им. Н.~Э.~Баумана)}
        }
    \end{minipage}

    \vspace{0.5cm}
    \rule{\linewidth}{0.4pt}

    % \begin{flushleft}
    %     {ФАКУЛЬТЕТ \uline{ИУ <<Информатика, искусственный интеллект и системы управления>> \hfill}}

    %     \vspace{0.5cm}

    %     {КАФЕДРА \uline{ИУ7 <<Программное обеспечение ЭВМ и информационные технологии>> \hfill}}
    % \end{flushleft}
\end{center}

\fontsize{11pt}{11pt}\selectfont
\begin{flushright}
    \begin{minipage}{0.4\textwidth}\raggedleft

        УТВЕРЖДАЮ \hspace{2.5cm}

        Заведующий кафедрой~\ulinetext[2cm]{(Индекс)}{ИУ7}

        \vspace{0.3cm}

        \ulinetext{}{} \ulinetext{\hfill (И. О. Фамилия)}{И. В. Рудаков}

        \vspace{0.3cm}

        << \ulinetext[1cm]{}{} >> \ulinetext{}{} \the\year \ г.
    \end{minipage}
\end{flushright}

\begin{center}\linespread{1}\selectfont
    \Large{\textbf{ЗАДАНИЕ}}

    \large{\textbf{на выполнение курсовой работы}}
\end{center}

\noindent по дисциплине \uline{\hspace{1cm} Компьютерная графика \hfill}

\noindent Студент группы \ulinetext{}{ИУ7-54Б}

\vspace{0.3cm}

\noindent \ulinetext[\textwidth]{(Фамилия, имя, отчество)}{Рунов Константин Алексеевич}

\vspace{0.3cm}

\noindent \begin{tabularx}{\linewidth}{@{}X@{}}
Тема курсовой работы \uline{
    Разработка программы для моделирования столкновений объектов в виртуальном пространстве.
    \hfill
    }
\end{tabularx}

% \noindent \uline{\hfill}

\vspace{0.3cm}

\noindent Направленность КР (учебная, исследовательская, практическая, производственная, др.)
\noindent \uline{\hfill учебная \hfill}

\noindent Источник тематики (кафедра, предприятие, НИР)
\noindent \ulinetext[8.3cm]{}{\hfill каферда \hfill}

\vspace{0.3cm}

\noindent График выполнения работы: 25\% к \ulinetext[0.5cm]{}{4} нед., 50\% к \ulinetext[0.5cm]{}{7} нед., 75\% к \ulinetext[0.5cm]{}{11} нед., 100\% к \ulinetext[0.5cm]{}{14} нед.

\vspace{0.3cm}

% @{} to eliminate edges of tabular
\noindent \begin{tabularx}{\linewidth}{@{}X@{}}
    \textbf{\textit{Задание}} \uline{
Разработать программу для моделирования столкновений объектов друг с другом. Пользователь должен иметь возможность создавать объекты: куб, шар, чайник, прямоугольный параллелепипед. Пользователь должен иметь возможность изменять свойста создаваемых объектов: размер, местоположение, цвет, масса, начальная скорость, а также иметь возможность изменять положение камеры и устанавливать значение гравитации. Программа должна моделировать реалистичное движение и столкновение объектов друг с другом, учитывая освещение сцены.
        \hfill}
\end{tabularx}

\vspace{0.3cm}

\noindent \textbf{\textit{Оформление курсовой работы:}}

\noindent Расчётно-пояснительная записка на \uline{25-30} листах формата А4.

% \noindent \uline{\hfill}

% \noindent \uline{\hfill}

\noindent \begin{tabularx}{\linewidth}{@{}X@{}}
    \uline{
Расчетно-пояснительная записка должна содержать постановку, введение, аналитическую часть, конструкторскую часть, технологическую часть, экспериментально-исследовательский раздел, заключение, список литературы, приложения.\hfill
        \hfill}
\end{tabularx}

\vspace{0.3cm}

\noindent Дата выдачи задания << \uline{\hspace{0.5cm}} >> \uline{\hspace{2cm}} \the\year \ г.

\vspace{0.3cm}

\noindent \textbf{Руководитель курсовой работы} \hfill \ulinetext{(Подпись, дата)}{} \ulinetext[4cm]{\hfill (И. О. Фамилия)}{\hfill А. А. Павельев}

\vspace{0.3cm}

\thispagestyle{empty}

\noindent \textbf{Студент} \hfill \ulinetext{(Подпись, дата)}{} \ulinetext[4cm]{\hfill (И. О. Фамилия)}{\hfill К. А. Рунов}

\vfill
